\documentclass[10pt,a4paper]{article}
\usepackage{fontspec,amsmath,amssymb,graphicx}
\usepackage[margin=18mm]{geometry}
\setmainfont{Times New Roman}
\setlength{\parindent}{0pt}
\setlength{\parskip}{5pt}
\allowdisplaybreaks
\sloppy
\emergencystretch=2em
\begin{document}
\section*{BT\_5\_4\_3\_01 — Lời giải bổ sung}
(a)--(c) Với hàm xuyên tâm $r^p$, $\Delta r^p=p(p+n-2)r^{p-2}$. Vì $r>0$, Laplace bằng 0 đúng khi $p=0$ hoặc $p=2-n$. Trong $n=3,p=-1$ cho $1/r$ điều hòa ngoài gốc.

(d) Đặt $v=x-\xi$, $|\xi|=1$, $1-|x|^2=-|v|^2-2v\cdot\xi$. Khi đó nhân Poisson bằng $-|v|^{2-n}-2(v\cdot\xi)|v|^{-n}$. Với $n\ne2$, hạng sau là đạo hàm theo hướng của $|v|^{2-n}$ chia $2-n$; do đó điều hòa. Với $n=2$, dùng $(v\cdot\xi)|v|^{-2}=\partial_\xi\ln|v|$, cũng điều hòa. Miền đúng là $\mathbb R^n\setminus\{\xi\}$; ký hiệu $\mathbb R^3$ ở dòng cuối đề chỉ phù hợp khi $n=3$.\end{document}